% Physiological Measurement submission, IOP Publishing template (iopjournal.cls, 2025 release,
% from publishingsupport.iopscience.iop.org). References: Harvard (author-year), as required by
% the journal. Supplementary material is compiled separately from supplementary_iop.tex.
\documentclass{iopjournal}

\usepackage{amssymb}
\usepackage{amsmath}
\usepackage{booktabs}
\usepackage{array}
\usepackage{multirow}
\usepackage{adjustbox}
\usepackage{makecell}
\usepackage{microtype}
\usepackage[round]{natbib}
\usepackage[section]{placeins}
\hypersetup{hidelinks}
\renewcommand{\topfraction}{0.9}
\renewcommand{\textfraction}{0.07}
\renewcommand{\floatpagefraction}{0.8}
\setcounter{topnumber}{3}

% running head without the template's placeholder journal details
\fancyhead[L]{\small Subject-independent surface EMG electrode reduction}
\fancyhead[C]{}
\fancyhead[R]{\small Ta \textit{et al}}

% kept so that generated tables compile; renders plain text
\newcommand{\rev}[1]{#1}

\input{generated/results_macros.tex}

\begin{document}

\title{Subject-independent surface EMG electrode reduction using cross-fitted attribution: a prespecified non-inferiority study}

\author{Minh-Tri Ta$^1$, Hoang-Khang Phan$^{2,3}$ and Le-Giang Tran$^{4,*}$}

\affil{$^1$College of Engineering and Computer Science, VinUniversity, Hanoi, Vietnam}

\affil{$^2$Department of Biomedical Engineering, Faculty of Applied Science, Ho Chi Minh City University of Technology (HCMUT), 268 Ly Thuong Kiet Street, District 10, Ho Chi Minh City, Vietnam}

\affil{$^3$Vietnam National University Ho Chi Minh City, Linh Trung Ward, Thu Duc City, Ho Chi Minh City, Vietnam}

\affil{$^4$Faculty of Electronics Technology, Industrial University of Ho Chi Minh City, Ho Chi Minh City, Vietnam}

\affil{$^*$Author to whom any correspondence should be addressed.}

\email{tranlegiang@iuh.edu.vn}

\keywords{surface electromyography, gesture recognition, electrode selection, SHAP, cross-fitting, non-inferiority, leave-one-subject-out}

\begin{abstract}
\textit{Objective.} To test, under prespecified decision rules, whether electrode selection guided by model attribution can reduce a forearm surface electromyography (EMG) array without meaningful loss of subject-independent gesture recognition.
\textit{Approach.} We analysed 43 GRABMyo participants (16 monopolar forearm electrodes, 16 active gestures) with leave-one-subject-out evaluation. Electrodes were ranked by TreeSHAP attributions cross-fitted over participant-disjoint training groups, with and without a stability penalty. Retrained subsets of top-ranked electrodes were compared with the full array against a prespecified non-inferiority margin of 0.02 Macro-F1, using participant-level cluster-bootstrap confidence intervals (CI).
\textit{Main results.} The full array reached a Macro-F1 of \FullXGB\ (95\% CI \FullXGBLCB--\FullXGBUCB). The smallest non-inferior subset contained \KStar\ electrodes (mean difference $\DeltaKTwelve$, 95\% CI $\DeltaKTwelveLCB$ to $\DeltaKTwelveUCB$), unchanged after Holm correction; eleven electrodes were not non-inferior. Attribution-ranked subsets exceeded the mean of random subsets at 4, 6, 8 and 12 electrodes, and deletion with retraining supported the faithfulness of the ranking. The stability penalty produced no measurable improvement. In exploratory comparisons, no statistically detectable difference was observed between attribution and an evenly spaced, data-free layout at 4, 8 and 12 electrodes. In an exploratory replication on putEMG (44 participants), a subset of \PutKStarPtTwo\ of 24 electrodes was non-inferior (\PutKStarHolmPtTwo\ with Holm correction).
\textit{Significance.} For offline, subject-independent recognition with this two-ring monopolar montage, about a quarter of the electrodes could be removed within a prespecified engineering tolerance, but not more; the advantage of attribution over even spatial coverage was detectable only at six electrodes.
\end{abstract}

\section{Introduction}
\label{sec:introduction}

Surface electromyography (EMG) pattern recognition is a widely used approach to multifunction myoelectric control \citep{hudgins1993new,englehart2003robust,scheme2011electromyogram,farina2014extraction}. Recognition generally improves with more recording sites, but every electrode adds cost, wiring, contact area and a point of failure in a wearable device \citep{vujaklija2016new}. This makes it attractive to identify a reduced set of electrodes that preserves recognition performance \citep{phinyomark2012feature,liu2014emg,altimemy2013classification,stango2015spatial,qu2021reduce}.

Three methodological problems limit what published electrode-reduction results can tell a device designer. First, most reductions are selected and evaluated within the same users. User-dependent evaluation is appropriate for personalised systems, but a device intended for many users faces a different question: a layout chosen for one person need not suit the next, and subject-independent recognition is generally harder than user-dependent recognition \citep{phinyomark2018emg,cote2019deep}. In addition, if data from a test user influence scaling, ranking or tuning, the reported subset is optimistically biased. Second, feature-attribution methods such as Shapley additive explanations \citep{lundberg2017unified,lundberg2020local} are increasingly used to rank channels \citep{salmi2025shapaad}, but attributions can be unstable across data splits \citep{visani2022statistical,kelodjou2023shaping,xiang2026roshap}, and a ranking computed on the evaluation data is itself a form of leakage. Third, the number of electrodes to retain is usually read off an accuracy curve by eye, with no prespecified tolerance and no statement of what `no meaningful loss' means \citep{schuirmann1987comparison,walker2011understanding}. A margin chosen after the curve is seen cannot support a non-inferiority claim \citep{campbell2018noninferiority}.

This study tests explainability-guided electrode reduction under conditions designed to remove these three sources of optimistic bias. Each held-out participant is excluded from feature standardisation, model fitting and electrode ranking. Rankings are computed by cross-fitting: TreeSHAP attributions are always evaluated on participants whom the attributing model did not see. The primary decision rules were fixed in a statistical analysis plan before any result with a reduced electrode set was available: the non-inferiority margin ($\delta = 0.02$ in Macro-F1), the participant as the statistical unit, the bootstrap, Holm correction within each confirmatory family of tests, and the definition of the smallest acceptable electrode count $k^*$. The detailed tests for the secondary hypotheses were specified later, and all departures from the plan are reported. We also test a stability-regularised cross-fitted attribution score (SR-CFA), which penalises electrodes whose attribution varies across disjoint groups of training participants.

The study is primarily a test, not a method proposal. Using the 43-participant GRABMyo dataset \citep{pradhan2022multiday}, we asked four prespecified questions:
(H1) is some subset of $k < 16$ electrodes non-inferior to the full array?
(H2) do attribution-ranked subsets outperform random subsets of the same size?
(H3) does the stability penalty improve on attribution magnitude alone?
(H4) is the attribution faithful, that is, does removing highly ranked electrodes reduce performance more than removing random or low-ranked ones?
We report every answer, including the negative ones. Exploratory analyses covered transfer across recording sessions, a bipolar re-referencing of the montage, and methodological replications on NinaPro DB5 \citep{pizzolato2017comparison} and putEMG \citep{kaczmarek2019putemg}, a 44-participant, 24-electrode dataset whose analysis was not part of the original plan. Figure~\ref{fig:overview} gives an overview of the study.

\begin{figure}[t]
\centering
\includegraphics[width=\textwidth]{figures/fig1_overview.pdf}
\caption{Overview of the study. (A) GRABMyo data and preprocessing: two circumferential rings of eight monopolar electrodes 2~cm apart on the forearm, 250~ms analysis windows and ten features per channel (160 features per window). The muscle drawings show general forearm anatomy only; electrodes were not registered to individual muscles. (B) Subject-independent validation: each participant is held out in turn (outer leave-one-subject-out split), and attributions are computed by cross-fitting over five participant-disjoint groups of the remaining participants. (C) The two electrode rankings compared in this study: attribution magnitude alone ($\lambda = 0$) and the stability-regularised score $Q_c = A_c - \lambda U_c$. The panel shows how the rankings are formed, not which performed better.}
\label{fig:overview}
\end{figure}

\section{Related work}
\label{sec:related_work}

\paragraph{Electrode reduction for myoelectric control}
Channel and feature reduction for surface EMG pattern recognition has been approached with feature-ranking and wrapper methods \citep{phinyomark2012feature,liu2014emg}, with channel selection on multi-channel and high-density arrays \citep{altimemy2013classification,stango2015spatial}, and with sparse-representation methods \citep{qu2021reduce}. Large-scale reviews of the field note that most evaluations are user-dependent and that cross-user models behave differently from per-user models \citep{phinyomark2018emg}. Deep cross-user models and transfer learning are one response to this gap \citep{cote2019deep}. The present study takes a different position: a conventional feature-based classifier, a strictly subject-independent protocol, and a question about hardware rather than model capacity.

\paragraph{Attribution-guided selection and attribution stability}
TreeSHAP computes Shapley-value attributions for tree ensembles \citep{lundberg2017unified,lundberg2020local}. Attribution has been used to reduce channels in other biosignal tasks, for example electroencephalography-based auditory attention decoding \citep{salmi2025shapaad}. Several studies show that attributions can change across resampling and neighbourhood choices \citep{visani2022statistical,kelodjou2023shaping}. A recent preprint treats Shapley attributions as distributions across resamples and proposes a robustness-adjusted importance score \citep{xiang2026roshap}. That idea is close to the stability penalty tested here, and we do not claim the penalty itself as new. The general idea of preferring features that are selected consistently across subsamples goes back to stability selection \citep{meinshausen2010stability}. The contribution of this study is the evaluation design: attribution is cross-fitted over participants inside each outer training fold, the stability penalty is tested as a prespecified hypothesis against magnitude alone, and the selected subsets are judged by a prespecified non-inferiority rule.

\paragraph{Faithfulness of attributions}
A ranking that looks plausible need not reflect what the model uses. Deletion-based tests remove the features that an explanation marks as important and measure the effect \citep{samek2017evaluating}. Retraining after removal avoids confounding the explanation with the distribution shift seen by a model that is not retrained \citep{hooker2019benchmark}. We use the retraining variant with electrodes as the unit.

\paragraph{Electrode shift and multi-day use}
Electrode displacement and day-to-day variation degrade myoelectric pattern recognition \citep{hargrove2006effect,kaufmann2010fluctuating}, and adaptation methods have been proposed to counter the resulting covariate shift \citep{vidovic2016improving,muceli2014extracting}. GRABMyo was recorded on three days \citep{pradhan2022multiday}, which allows an exploratory cross-session analysis. Because the evaluation here is subject-independent throughout, its cross-session contrasts combine a session effect with the cross-user gap and are not a pure test of donning and doffing.

\section{Methods}
\label{sec:methods}

Figure~\ref{fig:workflow} summarises the analysis workflow, which was applied separately to each dataset.

\begin{figure}[t]
\centering
\includegraphics[width=\textwidth]{figures/fig2_workflow.pdf}
\caption{Study workflow. For each held-out participant $p$, feature standardisation, model fitting and electrode ranking use only the remaining participants $R_{-p}$. Attributions are computed on participant-disjoint inner folds of $R_{-p}$, and models restricted to the top-ranked $k$ electrodes are retrained on $R_{-p}$ and tested on $p$. The same workflow was applied separately to each dataset.}
\label{fig:workflow}
\end{figure}

\subsection{Terminology}
An \emph{electrode} is a physical recording site, and a \emph{channel} is the signal it provides to the classifier; in the monopolar analyses each electrode yields one channel, so selecting electrodes and selecting channels are the same operation. The \emph{montage} is the spatial arrangement of the electrodes and the \emph{array} is the complete set. A \emph{bipolar pair} is the difference between two electrodes, derived from the monopolar recordings in one exploratory analysis. The subset size $k$ always counts the selection unit of the analysis in question (electrodes, or bipolar pairs).

\subsection{Datasets}

\paragraph{GRABMyo (primary)}
GRABMyo v1.1.0 \citep{pradhan2022multiday,jiang2024grabmyo}, distributed through PhysioNet \citep{goldberger2000physionet}, contains recordings from \GRABMyoParticipants\ able-bodied participants (\GRABMyoMale\ male, \GRABMyoFemale\ female; age \GRABMyoAgeMean\ $\pm$ \GRABMyoAgeSD\ years) on \GRABMyoSessions\ different days. Signals were recorded with an EMGUSB2+ amplifier (OT Bioelettronica) at 2048~Hz. In each session, participants performed 16 hand and wrist gestures and rest, with 7 trials of 5~s each. We analysed the 16 forearm electrodes. These form two circumferential rings of 8 electrodes (F1--F8 and F9--F16) 2~cm apart along the forearm, with electrode $j$ of one ring aligned with electrode $j+8$ of the other. Mean forearm circumference at the ring position was \GRABMyoCircMean\ $\pm$ \GRABMyoCircSD~cm, giving about \GRABMyoArcSpacing~cm between neighbouring electrodes in a ring. The recordings are monopolar, so the selection unit in the primary analysis is a single electrode. All \GRABMyoFilesVerified\ distributed files matched the dataset's published SHA-256 checksums, and all \GRABMyoRecords\ records were analysed. The primary task is the 16 active gestures (rest excluded), for which chance Macro-F1 is about 0.063.

\paragraph{Replication datasets}
Two further datasets were analysed separately, never pooled with GRABMyo, to examine whether the procedure behaves similarly with different hardware and tasks. Their channel numbers have no anatomical correspondence with those of GRABMyo, so only the behaviour of the procedure is compared across datasets.
NinaPro DB5 \citep{pizzolato2017comparison} provides \DBFiveSubjects\ able-bodied subjects (\DBFiveMale\ male, \DBFiveFemale\ female; age \DBFiveAgeMean\ $\pm$ \DBFiveAgeSD\ years) recorded with two Myo armbands (16 channels, 200~Hz). We used exercise~1 (12 finger movements, 6 repetitions each), segmented with the refined movement and repetition labels supplied with the dataset. The DB5 values are 8-bit integer codes, not physical units; their spectra show no powerline peak and extend to the Nyquist frequency.
putEMG \citep{kaczmarek2019putemg} provides 44 able-bodied participants (8 female, 36 male; aged 19--37 years) recorded on two days at least one week apart. Twenty-four monopolar electrodes form three circumferential bands of 8 (electrodes 1--8, 9--16 and 17--24) on the forearm, sampled at 5120~Hz with 12-bit resolution. Each session contains 20 repetitions of each of 7 active gestures, labelled from video-verified ground truth. The putEMG analysis was not part of the original plan; its design was fixed before any putEMG model was trained.

\paragraph{Ethics statement}
This study is a secondary analysis of previously published, de-identified datasets and involved no new participant recruitment or data collection. The original data collections were conducted in accordance with the Declaration of Helsinki and approved by the Office of Research Ethics of the University of Waterloo (ORE\#~31346; GRABMyo), the Ethics Commission of the Canton of Valais, Switzerland (NinaPro DB5) and the Bioethical Committee of Poznan University of Medical Sciences (no.~398/17; putEMG); informed-consent procedures are described in the dataset publications \citep{pradhan2022multiday,pizzolato2017comparison,kaczmarek2019putemg}. [AUTHOR INPUT REQUIRED: state whether an institutional review board determination or exemption was obtained for the present secondary analysis.]

\paragraph{Sex as a reported variable}
Sex was taken as reported by the source datasets (GRABMyo \GRABMyoMale\ male and \GRABMyoFemale\ female; DB5 \DBFiveMale\ male and \DBFiveFemale\ female; putEMG 36 male and 8 female). Sex-stratified analyses were not prespecified, and the study was not designed to estimate sex differences in recognition performance or electrode ranking; no sex-specific inference is made.

\subsection{Preprocessing and features}
\label{sec:preprocessing}
GRABMyo signals were band-pass filtered (20--450~Hz, 4th-order zero-phase Butterworth) and notch filtered at 50~Hz ($Q = 30$), the mains frequency identified from the spectra (mean band-power ratio 10.35 at 50~Hz and 0.62 at 60~Hz). Because the filters are zero-phase, the processing is non-causal and suitable only for offline evaluation. Windows were 250~ms with 50\% overlap and never crossed trial boundaries. From each channel we extracted ten features: mean absolute value (MAV), root mean square, variance, waveform length, zero crossings (ZC), slope sign changes (SSC), integrated EMG (IEMG), Willison amplitude (WAMP), mean frequency and median frequency. This gives a 160-dimensional vector per window: \GRABMyoWindows\ windows in total and \GRABMyoActiveWindows\ from active gestures.

The amplitude threshold of the ZC, SSC and WAMP features was set by one rule for every representation: 1.15 times the median absolute first difference of the signal, rounded to one significant figure. For GRABMyo it was computed once, from 54 unfiltered recordings of 6 participants, giving $10^{-2}$~mV; it was therefore not estimated within each outer fold. Recomputing it with each held-out participant excluded gave the same value in every fold (supplementary methods~S2), so fold-wise estimation would have produced identical features.
DB5 was band-pass filtered at 10--95~Hz without a notch filter, with 250~ms windows (50 samples) and a threshold of 7 quantisation steps from the same rule. putEMG was converted to millivolts and processed as GRABMyo (20--450~Hz band-pass, 50~Hz notch, 250~ms windows with 50\% overlap), using only windows that lay entirely within one gesture; because raw putEMG differences are quantised at one converter step, the rule was applied to the filtered signal and gave \PutThresh${}\times 10^{-3}$~mV, for 240 features per window. For both datasets the leave-one-out recomputation gave the same threshold in every fold.

\subsection{Classifiers and outer leave-one-subject-out evaluation}
All evaluation is leave-one-subject-out (LOSO), with one fold per participant. In fold $p$, standardisation statistics, model fitting and electrode ranking use only $R_{-p}$, and the test participant contributes only to the score. XGBoost was designated in the analysis plan as the primary classifier for the electrode-selection analyses because its computational cost made the required number of retrained models feasible; the plan was written before the full-cohort classifier comparison had finished, although results from an earlier, superseded run on 15 participants were available. XGBoost used 300 trees, depth 6, learning rate 0.1, row and column subsampling of 0.8, histogram-based tree construction and a multiclass softmax objective, and was trained on all \FoldTrainingWindows\ training windows of each fold. Two further classifiers were compared with all electrodes: a support vector machine with a radial basis function kernel ($C = 1$, kernel width set from the feature variance), trained on a stratified sample of 40,000 windows, and a random forest (200 trees, minimum leaf size 2) trained on 200,000 windows. These sample sizes were set by computational cost and chosen on two folds (supplementary methods~S1). Halving the support vector machine's training sample on one fold lowered Macro-F1 from \SVMCalibFortyK\ to \SVMCalibTwentyK, so its result is likely a lower bound. Hyperparameters were fixed in advance and not tuned. The primary metric is participant-level Macro-F1 over the 16 active gestures; accuracy and balanced accuracy were also recorded.

\subsection{Cross-fitted attribution and the SR-CFA score}
Within each outer fold, $R_{-p}$ was split into $K = 5$ participant-disjoint inner folds. For inner fold $j$, an XGBoost model was trained on the other inner folds, using a random sample of at most 200,000 windows standardised with inner-training statistics only, and exact TreeSHAP attributions \citep{lundberg2020local} were computed for up to 20,000 windows sampled at random from the inner-held-out participants. Because participants contribute similar numbers of windows (median \TestWindowsPerParticipant), each participant carries roughly equal weight. The importance of channel $c$ for window $i$ averaged the absolute attributions over the $G = 16$ classes and the $F = 10$ features of the channel,
\begin{equation}
\Phi_{i,c} = \frac{1}{G F}\sum_{g=1}^{G}\sum_{f=1}^{F} \lvert \phi_{i,g,c,f} \rvert ,
\end{equation}
and was averaged over windows and normalised to sum to one across channels, giving $I_{j,c}$. Across the five inner folds, $A_c = \operatorname{median}_j I_{j,c}$ measures attribution magnitude and $U_c = \operatorname{median}_j \lvert I_{j,c} - A_c \rvert$ its instability across disjoint groups of training participants. Electrodes were ranked by
\begin{equation}
Q_c(\lambda) = A_c - \lambda\, U_c , \qquad \lambda \in \{0, 0.25, 0.5, 1\},
\label{eq:srcfa}
\end{equation}
with ties broken by the lower channel number. The value $\lambda = 0$ gives cross-fitted SHAP magnitude and $\lambda > 0$ the stability-regularised score (SR-CFA). In every fold we verified that no inner-held-out participant contributed to the corresponding inner training set. The plan did not specify a criterion for choosing $\lambda$, so all four values are reported.

\subsection{Electrode-count performance curve and non-inferiority (H1)}
For every outer fold, $\lambda$ and $k \in \{1,\ldots,16\}$, XGBoost was retrained from scratch on the top-$k$ electrodes of that fold's ranking, using a fixed, stratified sample of 200,000 training windows of the fold, and tested on the held-out participant. The $k = 16$ model of the same curve is the reference, so each paired difference $\Delta_{p,k} = \mathrm{F1}_{p,k} - \mathrm{F1}_{p,16}$ compares models trained on identical data. Subset size $k$ is non-inferior at margin $\delta$ when $L_k$, the lower bound of the 95\% participant-cluster bootstrap confidence interval of the mean difference $\bar\Delta_k$, exceeds $-\delta$. The smallest non-inferior subset is $k^* = \min\{k < 16 : L_k > -\delta\}$, recorded as `none' if no $k$ qualifies. As required by the plan, the family of 15 tests was also corrected by the Holm procedure, using one-sided cluster-bootstrap $p$-values for $H_0\colon \bar\Delta_k \le -\delta$ at $\alpha = 0.025$.

The primary margin $\delta = 0.02$ is an absolute difference in Macro-F1 on a 0--1 scale. It was fixed before the curve was computed as the largest loss we would accept in exchange for fewer electrodes. Relative to the full-array reference it corresponds to about \MarginPctOfRef\% of performance and about \MarginPctAboveChance\% of the distance above chance (1/16). It is an engineering tolerance and has not been validated against clinical outcomes or user performance. Because $k^*$ depends strongly on $\delta$, the prespecified sensitivity margins $\delta = 0.01$ and $0.03$ are also reported. Non-inferiority within this tolerance does not imply equivalence.

\subsection{Comparison with random subsets and other selection methods (H2)}
At $k \in \{4, 6, 8, 12\}$ we compared six rankings, each computed on the training participants only:
\begin{itemize}
\item SR-CFA ($\lambda = 0.5$);
\item cross-fitted SHAP ($\lambda = 0$);
\item permutation importance: a model was fitted on 20,000 training windows and each of the 160 features was permuted in 8,000 further windows (3 repeats, Macro-F1 scoring); an electrode's score is the mean over its 10 features;
\item mutual information between each feature and the gesture label on the same 20,000 windows, averaged over the electrode's 10 features;
\item geometric selection: farthest-point sampling on the $2 \times 8$ montage (arc distance around a 24~cm circumference, 2~cm between rings, starting from F1, ties to the lower index), which uses no data;
\item random subsets: 100 independent random subsets per participant at each $k$, and 500 at the prespecified $k^*$ (seed fixed per participant).
\end{itemize}
Evaluating hundreds of random subsets per participant required a lighter comparison model (XGBoost with 50 trees, trained on a stratified sample of 150,000 windows), applied identically to every method. Its absolute scores are lower than those of the performance curve and are compared only with each other. The confirmatory H2 family contrasts each attribution ranking (SR-CFA and cross-fitted SHAP) with the participant's mean over its random subsets (Holm over 8 tests). Comparisons with permutation importance, mutual information and geometric selection are exploratory and uncorrected. We also report the percentile of the SR-CFA subset within each participant's random subsets. Below, `SR-CFA' refers to $\lambda = 0.5$ and `cross-fitted SHAP' to $\lambda = 0$.

\subsection{Stability regularisation (H3) and faithfulness (H4)}
H3 compared each $\lambda > 0$ with $\lambda = 0$ on the performance curve at $k \in \{4, 6, 8, 12\}$ (Holm over 12 tests). For H4 we used deletion with retraining \citep{hooker2019benchmark}. For $n \in \{2, 4, \ldots, 14\}$ electrodes retained, the comparison model was retrained after removing the $16 - n$ highest-ranked electrodes (Remove-Top), the $16 - n$ lowest-ranked electrodes (Remove-Bottom), or random electrodes (Remove-Random, the random subsets above, available for $n \in \{4, 6, 8, 12\}$). A faithful ranking gives Remove-Top $<$ Remove-Random $<$ Remove-Bottom. Paired contrasts were Holm-corrected within H4.

\subsection{Statistics}
The participant is the statistical unit ($n = 43$ for GRABMyo, 10 for DB5 and 44 for putEMG), and windows are never treated as independent. Confidence intervals (CI) are 95\% percentile intervals from 10,000 cluster-bootstrap resamples of participants (seed 20260728). Paired comparisons use the Wilcoxon signed-rank test with the matched-pairs rank-biserial correlation $r_{rb}$ as effect size. Holm correction is applied within each confirmatory family; exploratory analyses are uncorrected and labelled as such. A non-significant difference is reported as `no statistically detectable difference', never as equivalence, because no equivalence test was performed. No participant was excluded.

\subsection{Exploratory analyses}
\paragraph{Ranking stability} We compared rankings derived from single participants' attributions (all 903 pairs), from the five inner folds and from the 43 outer folds, using Spearman's $\rho$, Kendall's $\tau$ and the Jaccard index of the top-$k$ sets. A participant's attribution profile is the mean of their normalised attribution over the outer folds in which they were held out of an inner model. Outer folds share about 98\% of their training participants, so their agreement is only an upper bound.

\paragraph{Replication} The LOSO baseline, cross-fitted ranking ($K = 5$) and performance curve were repeated on DB5 and putEMG, each with its full array as the reference ($k = 16$ and $k = 24$), using all training windows and the same bootstrap, margins and $k^*$ rule, with Holm-corrected results also reported. On putEMG, ranked subsets were also compared with ten random subsets per participant at $k \in \{6, 12, 18\}$ (Wilcoxon, Holm over 3). No setting was tuned on either dataset.

\paragraph{Cross-session} XGBoost was trained on session~1 of $R_{-p}$ and tested on each session of $p$, and trained on sessions 1 and 2 and tested on session~3. The second condition also has about twice the training data, so it confounds session diversity with data quantity.

\paragraph{Bipolar re-referencing} To use a bipolar pair as the selection unit, we formed $B_j = F_j - F_{j+8}$ ($j = 1, \ldots, 8$) from the monopolar recordings, re-extracted features with the threshold given by the same rule ($4 \times 10^{-3}$~mV; the bipolar signal is about half the monopolar amplitude), and repeated the analysis with the $k = 8$ model as reference. These are derived pairs, not physical bipolar electrodes.

\subsection{Departures from the analysis plan}
\label{sec:deviations}
The detailed H2--H4 test families were specified after the underlying results were available; the plan's Holm correction for H1 was applied after the primary analysis; $\lambda$ was not selected; training-sample sizes were chosen on two folds; the putEMG replication, a fixed-layout check and a feature-redundancy check were added; and two exploratory analyses listed in the plan were not run. None of these changes altered a prespecified decision rule. The complete list is given in supplementary methods~S1.

\section{Results}
\label{sec:results}

\subsection{Full-array subject-independent performance}
With all 16 electrodes, XGBoost reached a Macro-F1 of \FullXGB\ (95\% CI \FullXGBLCB--\FullXGBUCB) over the 43 held-out participants. The support vector machine reached \FullSVM\ (\FullSVMLCB--\FullSVMUCB) and the random forest \FullRF\ (\FullRFLCB--\FullRFUCB) (table~\ref{tab:full_channel}, figure~\ref{fig:full}). No statistically detectable difference was observed between XGBoost and the support vector machine (paired difference $\XGBvsSVMDiff$, 95\% CI $\XGBvsSVMLCB$ to $\XGBvsSVMUCB$, $p_{\mathrm{Holm}} \XGBvsSVMP$, $r_{rb} = \XGBvsSVMR$). Both exceeded the random forest ($p_{\mathrm{Holm}} \XGBvsRFP$ and $\SVMvsRFP$; $r_{rb} = \XGBvsRFR$ and $\SVMvsRFR$). The spread between participants was far larger than any difference between models: XGBoost ranged from \FullXGBMin\ to \FullXGBMax\ across participants.

\input{tables/tab_full_channel.tex}

\begin{figure}[t]
\centering
\includegraphics[width=\textwidth]{figures/fig3_full_channel.pdf}
\caption{Full-array leave-one-subject-out performance on GRABMyo (43 participants). (a) Macro-F1 of three classifiers; each point is one held-out participant, grey lines join the same participant, black bars mark medians and diamonds mark means with 95\% cluster-bootstrap confidence intervals. The support vector machine and random forest were trained on samples of 40,000 and 200,000 windows. (b) Paired XGBoost minus support vector machine difference per participant, sorted; the grey band is the 95\% confidence interval of the mean difference.}
\label{fig:full}
\end{figure}

\subsection{H1: a 12-electrode subset is non-inferior within the prespecified margin}
Performance declined smoothly as electrodes were removed (figure~\ref{fig:curve}, table~\ref{tab:sparse_ni}). The reference of the curve ($k = 16$) scored \CurveKSixteen, slightly below the full-data result because all curve models share one 200,000-window training sample. At $\delta = 0.02$, the smallest non-inferior subset was $k^* = \KStar$ electrodes ($\bar\Delta_{12} = \DeltaKTwelve$, 95\% CI $\DeltaKTwelveLCB$ to $\DeltaKTwelveUCB$). At $k = 11$ the lower bound ($\DeltaKElevenLCB$) crossed the margin. Every $k$ from 12 to 15 was non-inferior, so the result is not an isolated crossing. With Holm correction over the 15 subset sizes, $k^*$ was unchanged at all three margins ($k = 12$: $p_{\mathrm{Holm}} \HolmPKTwelve$; supplementary table~S1). The sensitivity margins gave $k^* = \KStarMarginPtOne$ at $\delta = 0.01$ and $k^* = \KStarMarginPtThree$ at $\delta = 0.03$. The 12-electrode subset was therefore non-inferior within the prespecified engineering tolerance, a reduction of \KStarReductionPct\% of the array. Larger reductions were not: halving the array to $k = 8$ cost $\DeltaKEight$ (95\% CI $\DeltaKEightLCB$ to $\DeltaKEightUCB$), and $k = 4$ cost $\DeltaKFour$. All paired differences at $k < 16$ were negative. The electrodes most often excluded at $k = 12$ were \DroppedFourSet\ (\DroppedFourFolds\ of 43 folds).

\input{tables/tab_sparse_ni.tex}

\begin{figure}[t]
\centering
\includegraphics[width=\textwidth]{figures/fig4_sparse_curve.pdf}
\caption{Electrode-count performance curve on GRABMyo (43 participants). (a) Mean Macro-F1 of models retrained on the top-$k$ electrodes of each fold's cross-fitted ranking, for the four values of the stability weight $\lambda$ (the curves overlap); the band is the 95\% cluster-bootstrap confidence interval for $\lambda = 0$. (b) Paired difference to the $k = 16$ model with 95\% confidence interval. The thick dashed line marks the primary margin ($-\delta$, $\delta = 0.02$); the thin dotted and dash-dotted grey lines mark the sensitivity margins (0.01 and 0.03).}
\label{fig:curve}
\end{figure}

\subsection{H2: attribution-ranked subsets exceed the mean of random subsets}
Under the shared comparison model (table~\ref{tab:baselines}, figure~\ref{fig:baselines}), SR-CFA subsets outperformed the mean of random subsets at every tested $k$ (Holm-corrected, 8 tests). The advantage was largest at intermediate sizes: $\HTwoKSixDiff$ Macro-F1 at $k = 6$ (95\% CI $\HTwoKSixLCB$ to $\HTwoKSixUCB$, $p_{\mathrm{Holm}} \HTwoKSixP$, $r_{rb} = \HTwoKSixR$) and $\HTwoKEightDiff$ at $k = 8$. It was small at $k = 4$ ($\HTwoKFourDiff$) and at $k = 12$ ($\HTwoKTwelveDiff$, $p_{\mathrm{Holm}} \HTwoKTwelveP$). Within each participant's own random subsets, the SR-CFA subset had a median percentile of \PctKFour\ at $k = 4$, \PctKSix\ at $k = 6$, \PctKEight\ at $k = 8$ and \PctKTwelve\ at $k = 12$, and it reached the 95th percentile for \PctKAboveNinetyFiveSix\ of 43 participants at $k = 6$ and \PctKAboveNinetyFiveTwelve\ at $k = 12$. Attribution-ranked subsets therefore exceeded the average random subset at every size tested but seldom reached the upper tail of a participant's random subsets.

In exploratory, uncorrected contrasts, no statistically detectable difference between SR-CFA and permutation importance was observed at $k = 4$, 8 or 12 (for example, $k = 8$: $\VSPermKEightDiff$, $p \VSPermKEightP$); at $k = 6$, SR-CFA was higher by $\VSPermKSixDiff$ ($p \VSPermKSixP$). Likewise, no statistically detectable difference between SR-CFA and geometric selection, which spreads electrodes evenly around both rings without using data, was observed at $k = 4$ ($\VSGeomKFourDiff$, $p \VSGeomKFourP$), $k = 8$ ($\VSGeomKEightDiff$, $p \VSGeomKEightP$) or $k = 12$ ($\VSGeomKTwelveDiff$, $p \VSGeomKTwelveP$); at $k = 6$, SR-CFA was higher by $\VSGeomKSixDiff$ ($p \VSGeomKSixP$). Mutual information gave the lowest scores at every $k \le 8$.

\input{tables/tab_baselines.tex}

\begin{figure}[t]
\centering
\includegraphics[width=\textwidth]{figures/fig5_selection_baselines.pdf}
\caption{Selection methods at matched subset size $k$ (GRABMyo, 43 participants; one comparison model shared by all methods). (a) Mean Macro-F1 with 95\% cluster-bootstrap confidence intervals, grouped by method type: attribution rankings (stability-regularised cross-fitted attribution, SR-CFA, and cross-fitted SHAP), model-based rankings (permutation importance and mutual information), data-free geometric selection and the mean over random subsets. (b) Percentile of the SR-CFA subset within each participant's own random subsets (one point per participant; bar = median; 100 random subsets per participant, 500 at $k = 12$).}
\label{fig:baselines}
\end{figure}

\subsection{H3: the stability penalty produced no measurable improvement}
The instability term was small relative to attribution magnitude: the median $U_c/A_c$ was \UOverAMedian\% (maximum \UOverAMax\%). Penalising it therefore rarely changed a ranking. Relative to $\lambda = 0$, the top-12 set changed in \TopTwelveDiffersLambdaPtTwoFive, \TopTwelveDiffersLambdaPtFive\ and \TopTwelveDiffersLambdaOne\ of 43 folds at $\lambda = 0.25$, 0.5 and 1, and across all four $k$ tested at least \HThreeMinTied\ of 43 participants received the identical subset under both rankings. None of the \HThreeNTests\ contrasts was significant after Holm correction (largest mean difference \HThreeMaxAbsDiff; smallest $p_{\mathrm{Holm}} \HThreeMinP$; supplementary table~S2). The four curves in figure~\ref{fig:curve}(a) differ by at most \LambdaCurveMaxSpread, and $k^*$ was identical under every $\lambda$ and margin. H3 was not supported.

\subsection{H4: deletion with retraining supports faithfulness}
The expected ordering, Remove-Top $<$ Remove-Random $<$ Remove-Bottom, held at every number of electrodes retained (table~\ref{tab:faithfulness}, figure~\ref{fig:faithfulness}); \FaithNSig\ of the \FaithNTests\ Holm-corrected contrasts were significant. With six electrodes retained, Remove-Top scored \FaithTopKSix, Remove-Random \FaithRandomKSix\ and Remove-Bottom \FaithBottomKSix\ (Remove-Top minus Remove-Random $\FaithTRKSixDiff$, $p_{\mathrm{Holm}} \FaithTRKSixP$; Remove-Random minus Remove-Bottom $\FaithRBKSixDiff$, $p_{\mathrm{Holm}} \FaithRBKSixP$). The only non-significant contrast was Remove-Top versus Remove-Bottom with 14 electrodes retained ($\FaithTBKFourteenDiff$, $p_{\mathrm{Holm}} \FaithTBKFourteenP$), where only two electrodes are removed. The separation between deletion orders was largest at intermediate sizes and small at the extremes.

\input{tables/tab_faithfulness.tex}

\begin{figure}[t]
\centering
\includegraphics[width=0.62\textwidth]{figures/fig6_faithfulness.pdf}
\caption{Deletion-with-retraining test of attribution faithfulness on GRABMyo (43 participants; mean and 95\% cluster-bootstrap confidence interval). Remove-Top deletes the highest-ranked electrodes, Remove-Bottom the lowest-ranked and Remove-Random a random set; all use the same comparison model and training sample. Remove-Random is shown where random subsets were sampled (4, 6, 8 and 12 electrodes retained).}
\label{fig:faithfulness}
\end{figure}

\subsection{Stability of the ranking across participants}
Attribution was concentrated and reproducible (figure~\ref{fig:attribution}). The same four electrodes (\TopFourSet) formed the top four in \TopFourFolds\ of 43 outer folds and carried \TopFourShare\% of the summed attribution, whereas the eight lowest-ranked electrodes carried \BottomEightShare\%. Rankings computed from individual participants' attributions agreed closely (table~\ref{tab:stability}): mean Spearman $\rho$ = \StabSubjRho\ $\pm$ \StabSubjRhoSD\ over \StabSubjPairs\ pairs, and a mean Jaccard index of the top six of \StabSubjJacSix. Rankings from disjoint inner folds agreed at $\rho$ = \StabInnerRho. The outer-fold agreement ($\rho$ = \StabOuterRho) is an upper bound only, because those training sets share about 98\% of their participants. Across folds, the highest- and lowest-ranked electrodes kept their positions, and the rank ranges were widest in the middle of the ordering (figure~\ref{fig:attribution}(b)).

\input{tables/tab_stability.tex}

\begin{figure}[t]
\centering
\includegraphics[width=\textwidth]{figures/fig7_attribution.pdf}
\caption{Cross-fitted electrode attribution on GRABMyo ($\lambda = 0$, 43 outer folds). (a) Mean share of attribution per electrode, drawn on the $2 \times 8$ montage; each ring~2 electrode is aligned with, and 2~cm from, the ring~1 electrode at the same circumferential position. (b) Rank of each electrode across folds (median and range).}
\label{fig:attribution}
\end{figure}

\subsection{Replication on independent datasets}
\paragraph{putEMG (exploratory)} With all 24 electrodes, leave-one-subject-out Macro-F1 over the 44 participants was \PutFull\ (95\% CI \PutFullLCB--\PutFullUCB; range \PutFullMin--\PutFullMax; 7 gestures, chance about 0.143) (figure~\ref{fig:generalisation}(d), supplementary table~S3). At $\delta = 0.02$ the smallest non-inferior subset was $k^* = \PutKStarPtTwo$ of 24 electrodes ($\bar\Delta = \PutDeltaKFifteen$, 95\% CI $\PutDeltaKFifteenLCB$ to $\PutDeltaKFifteenUCB$), while at $k = \PutKStarMinusOne$ the lower bound was $\PutDeltaKFourteenLCB$; every larger subset was also non-inferior, and the decision at $k^*$ held under \PutSeedFrac\% of 200 alternative bootstrap seeds. With Holm correction over the 23 subset sizes, the smallest non-inferior subset was $\PutKStarHolmPtTwo$ electrodes. The sensitivity margins gave $k^* = \PutKStarPtOne$ at $\delta = 0.01$ and $\PutKStarPtThree$ at $\delta = 0.03$. Ranked subsets exceeded the mean of ten random subsets at all three sizes tested ($k = 6$: $\PutVsRandKSixDiff$, 95\% CI $\PutVsRandKSixLCB$ to $\PutVsRandKSixUCB$; $k = 12$: $\PutVsRandKTwelveDiff$; $k = 18$: $\PutVsRandKEighteenDiff$; each $p_{\mathrm{Holm}} \PutVsRandKSixP$). The ranking was stable, with the same top four electrodes (\PutTopFourSet) in \PutTopFourFolds\ of 44 folds. Because putEMG differs in montage, electrode count and gesture set, this is a methodological replication: the behaviour of the procedure, not the retained fraction, is comparable with GRABMyo.

\paragraph{NinaPro DB5 (exploratory)} With all 16 channels, Macro-F1 was \DBFiveFull\ $\pm$ \DBFiveFullSD\ (12 movements; chance about 0.083), far below GRABMyo and with wide between-subject spread (\DBFiveFullMin--\DBFiveFullMax) (figure~\ref{fig:generalisation}(b), supplementary table~S4). At $\delta = 0.02$, $k = 12$ met the margin by a negligible amount ($\bar\Delta_{12} = \DBFiveDeltaKTwelve$, lower bound $\DBFiveDeltaKTwelveLCBSix$): it was non-inferior under \DBFiveSeedFrac\% of \DBFiveSeedN\ alternative bootstrap seeds, $k = 13$, 14 and 15 all failed (lower bounds $\DBFiveDeltaKThirteenLCB$, $\DBFiveDeltaKFourteenLCB$ and $\DBFiveDeltaKFifteenLCB$), and with Holm correction no subset was non-inferior. With ten subjects the estimates have wide uncertainty, and DB5 neither supports nor contradicts the GRABMyo reduction; its mean curve increased broadly with $k$ but was not monotone.

\begin{figure}[t]
\centering
\includegraphics[width=\textwidth]{figures/fig8_generalisation.pdf}
\caption{Exploratory analyses. (a) GRABMyo cross-session evaluation: each point is one held-out participant, grey lines join the same participant across conditions and bars mark medians; every condition is also subject-independent. (b--d) Mean paired difference to the full array of the same curve, $\bar\Delta_k$, with 95\% cluster-bootstrap confidence interval, for (b) NinaPro DB5 (10 subjects, reference $k = 16$), (c) GRABMyo derived bipolar pairs (43 participants, reference $k = 8$) and (d) putEMG (44 participants, reference $k = 24$). The dashed line marks the primary margin; filled markers indicate a lower bound above the margin and open markers a lower bound at or below it.}
\label{fig:generalisation}
\end{figure}

\subsection{Other exploratory analyses}
\paragraph{Cross-session} Training on session~1 and testing on sessions 1, 2 and 3 gave \CSOneSOne, \CSOneSTwo\ and \CSOneSThree\ (supplementary table~S6, figure~\ref{fig:generalisation}(a)). Mean Macro-F1 on sessions 2 and 3 was numerically higher than on session 1; the session-2 difference was nominally significant ($\CSSTwoVsSOneDiff$, 95\% CI $\CSSTwoVsSOneLCB$ to $\CSSTwoVsSOneUCB$, uncorrected $p \CSSTwoVsSOneP$), and no statistically detectable difference was observed for session 3 ($\CSSThreeVsSOneDiff$, $p \CSSThreeVsSOneP$). Because every condition is also subject-independent, this does not show that sessions are interchangeable, only that no additional cross-day loss was detectable on top of the cross-user gap. Training on sessions 1 and 2 improved session-3 performance over training on session 1 alone ($\CSTwoVsOneDiff$, 95\% CI $\CSTwoVsOneLCB$ to $\CSTwoVsOneUCB$, $p \CSTwoVsOneP$), but this design cannot separate day diversity from the doubled amount of training data.

%BIPOLAR_RESULTS%

\section{Discussion}
\label{sec:discussion}

\paragraph{How far the array can be reduced}
Under a margin fixed before the curve was computed, four of the 16 GRABMyo electrodes could be removed while remaining non-inferior to the full array; five could not. Every $k$ from 12 to 15 met the margin, so the result does not rest on a single isolated crossing, and it was unchanged by Holm correction. The margin is an engineering tolerance, not a clinical threshold, and every reduction cost some Macro-F1: halving the array cost about \DeltaKEight. In a subject-independent setting this 16-electrode monopolar array therefore contains only limited redundancy for 16-gesture recognition. User-dependent studies that retain far fewer channels answer a different question \citep{phinyomark2018emg,qu2021reduce}, and their results should not be taken as evidence that a layout intended for many users can be equally sparse.

\paragraph{What attribution contributes}
Attribution-ranked subsets exceeded the mean of random subsets at every size tested (H2), and deletion with retraining showed the expected ordering at every size (H4), which supports that the ranking reflects electrodes on which the retrained models depend. The mean random subset is, however, a limited benchmark. In exploratory comparisons, no statistically detectable difference was observed between the attribution ranking and either permutation importance or data-free geometric selection at 4, 8 and 12 electrodes; at 6 electrodes, the attribution ranking was higher by $\VSGeomKSixDiff$ than geometric selection and by $\VSPermKSixDiff$ than permutation importance. Within each participant's own random subsets, the attribution subset was typically above the median but rarely near the top. One explanation consistent with this pattern is that even circumferential coverage accounts for much of what any ranking achieves. Absence of a detectable difference is not evidence of equivalence, and whether an evenly spaced layout is itself non-inferior to the full array was not tested.

\paragraph{What the stability penalty contributes}
The stability penalty produced no measurable improvement under this protocol (H3). With cross-fitting over roughly 34 participants per inner training set, inner-fold attributions agreed so closely that the instability term was a few percent of the magnitude term, and the penalty changed the top-12 set in at most one fold. In these data, plain cross-fitted TreeSHAP magnitude is therefore the simpler option. This does not show that stability regularisation is ineffective in general; it may matter with fewer training participants or less stable attributions, and a related robustness-weighted score has been proposed independently \citep{xiang2026roshap}.

\paragraph{What the replications show}
On putEMG, with 44 participants and a 24-electrode montage in three bands, the same procedure again produced a stable ranking that outperformed random subsets and a non-inferior reduction to $\PutKStarPtTwo$ of 24 electrodes ($\PutKStarHolmPtTwo$ with Holm correction), with no larger subset failing the margin. This analysis was not in the original plan and is a methodological replication on a distinct dataset and montage, not evidence that a particular retained fraction generalises across hardware. On DB5, ten subjects and a low baseline left the estimates with wide uncertainty: the decision at $k = 12$ depended on the bootstrap seed and did not survive Holm correction, so DB5 neither supports nor contradicts the GRABMyo reduction. The primary conclusions do not depend on DB5.

\paragraph{Positional and physiological interpretation}
The four highest-ranked electrodes (\TopFourSet) were the same in \TopFourFolds\ of 43 folds and in single-participant rankings (mean Spearman $\rho$ = \StabSubjRho). We deliberately keep the interpretation positional. The dataset documents electrode positions more reliably than the muscles beneath them, no ultrasound or other imaging-based registration of electrodes to muscles was available, and surface EMG at each site reflects volume-conducted activity from several muscles \citep{farina2004extraction}. High attribution concentrated at a few circumferential positions, the top and bottom of the ranking were stable, and the ordering varied mainly among mid-ranked electrodes. The spatial ranking is compatible with a physiologically plausible regional structure, in which a few sites capture activity shared by several gestures, but channel-to-muscle correspondence was not independently verified.

\paragraph{Sessions and sensor type}
In the exploratory cross-session analysis, no additional loss on later days was detected on top of the cross-user gap. This does not show that electrode shift is harmless \citep{hargrove2006effect,kaufmann2010fluctuating,vidovic2016improving}, and the gain from training on two sessions is confounded with doubling the training data. The primary unit here is a monopolar electrode, whereas conventional devices measure bipolar pairs, and removing four electrodes does not remove four sensors of such a device. %BIPOLAR_DISCUSSION% The 25\% reduction is therefore a statement about monopolar electrodes on this montage, not evidence of an equivalent reduction for a physical bipolar armband; physical bipolar electrodes may also behave differently from derived pairs.

\paragraph{Limitations}
All participants were able-bodied, and all analysis was offline: filtering was zero-phase (non-causal) and windows were 250~ms, so the results do not establish performance in real-time control. Fatigue, limb posture and electrode shift were not isolated. No imaging-based registration of electrodes to muscles was available. Rankings were computed from full-array models and could differ if models were trained only on a reduced array. The support vector machine was trained on a sample of windows for computational reasons and is likely a lower bound, and the H2 and H4 comparisons used a lighter model than the performance curve. The detailed H2--H4 tests and the Holm correction for H1 were specified after the primary results were available, $\lambda$ was not selected, and training-sample sizes were chosen on two folds. The putEMG replication and the fixed-layout and feature-redundancy checks were added after the plan, and the fixed-layout result (supplementary methods~S4) is a single-split consistency check only. DB5 had only ten subjects. The bipolar analysis used pairs derived from monopolar recordings, which may behave differently from physical bipolar electrodes. The window-length and 17-class analyses listed in the plan were not run.

\section{Conclusion}
\label{sec:conclusion}

For offline, subject-independent gesture recognition with a two-ring monopolar forearm montage in able-bodied participants, four of 16 electrodes could be removed while remaining non-inferior to the full array within a prespecified 0.02 Macro-F1 engineering tolerance, but five could not. Cross-fitted attribution produced rankings that outperformed the mean of random subsets and whose faithfulness was supported by deletion with retraining, whereas the stability penalty produced no measurable improvement. In an exploratory replication on a larger, differently configured dataset (putEMG), a stable ranking again outperformed random subsets and \PutKStarPtTwo\ of 24 electrodes met the same tolerance.

The study does not establish clinical performance, a reduction of physical bipolar hardware, an advantage of attribution over simple evenly spaced layouts, real-time performance, muscle-level localisation, or generalisation beyond these cohorts and montages. Validation in target clinical populations, such as people with upper-limb loss, and in online, causal control settings is required before deployment claims can be made.


\ack{[AUTHOR INPUT REQUIRED: acknowledgements, if any.]}

\funding{[AUTHOR INPUT REQUIRED: funding source and grant number, or `This research received no specific grant from any funding agency in the public, commercial or not-for-profit sectors.' only if true.]}

\roles{[AUTHOR INPUT REQUIRED: contributions of Minh-Tri Ta, Hoang-Khang Phan and Le-Giang Tran using CRediT roles.]}

\section*{Conflict of interest}
The authors declare no conflicts of interest.

\data{GRABMyo v1.1.0 is available from PhysioNet \citep{jiang2024grabmyo} (\url{https://doi.org/10.13026/89dm-f662}); NinaPro DB5 from the NinaPro repository (\url{https://ninapro.hevs.ch}); and putEMG from its authors under a CC BY-NC 4.0 licence (\url{https://www.biolab.put.poznan.pl/putemg-dataset/}). The statistical analysis plan, the record of analysis decisions, the analysis code, the environment specification and all result files from which the numbers in this paper are computed are archived at [AUTHOR INPUT REQUIRED: INSERT PERMANENT REPOSITORY DOI].}

\suppdata{Supplementary methods S1--S5 (departures from the analysis plan, threshold checks, deletion-analysis consistency, fixed-layout and feature-redundancy analyses) and supplementary tables S1--S6 are provided in a separate file.}

\bibliographystyle{jphysicsB}
\bibliography{references}

\end{document}
